\chapter{关键结果汇总与结论}

\section{关键结果汇总}

表\ref{tab:final_results}汇总本文五个问题的最终输出。该表可作为全文结果索引，也便于检查题目要求是否全部覆盖。

\begin{table}[htbp]
\centering
\footnotesize
\caption{五个问题关键结果汇总}
\label{tab:final_results}
\begin{tabularx}{0.98\textwidth}{C{1.2cm}X X}
\toprule
问题 & 最终模型 & 关键结果 \\
\midrule
1 & Huber稳健仿射校准 & $\hat y=-0.5264+0.8817x$；待校正值为5.762、15.809、73.836、108.415、147.311 mm \\
2 & Hampel速度清洗+复合证据变点 & 阶段转换主节点为第8108号和第9501号；三阶段平均速度为0.262、1.859、6.234 mm/h \\
3 & 状态空间补齐+变量专属异常识别 & 单变量异常63个、共同异常17个；贡献排序为孔隙水压力、深部位移、降雨量、微震事件数 \\
4 & 阶段归一化基线+GBDT残差 & 指定时刻预测值为5.163、30.767、60.259、383.185、383.697 mm，并给出95\%预测区间 \\
5 & 五变量组合验证+三级速度预警 & 保留降雨量、孔隙水压力、微震事件数、爆破点距离、单段最大药量；三级阈值为2.980、6.418、9.949 mm/h \\
\bottomrule
\end{tabularx}
\end{table}

\section{工程预警实施流程}

若将本文模型用于实际边坡监测，可按图式流程执行：

\begin{enumerate}
  \item 每10 min更新一次位移、降雨、孔压、微震和爆破记录，并按本文规则处理爆破空值。
  \item 对位移差分速度进行Hampel清洗，计算6 h滚动中位速度。
  \item 根据阶段归一化基线判断当前处于缓慢、加速或快速形变阶段，并计算扰动残差。
  \item 当速度超过关注阈值且持续不少于2 h时，进入关注状态；若同时出现逆速度下降和扰动增强，则升级为预警。
  \item 当速度超过严重预警阈值并持续不少于6 h，或逆速度外推失稳时间进入24 h窗口时，进入严重预警状态。
\end{enumerate}

该流程把“趋势识别”和“扰动解释”分开：趋势识别负责判断坡体是否进入危险演化阶段，扰动解释负责判断风险是否由降雨、孔压、微震或爆破进一步放大。

\section{结论}

本文围绕边坡预警问题建立了完整的数学建模方案。问题一证明两类位移传感器可用稳健线性关系统一，避免了高阶校准的不稳定性。问题二表明阶段转换不能只看累计位移或单点跳变，必须综合速度水平、位移曲率和持续加速度。问题三说明多源监测变量具有不同统计性质，变量专属补齐和异常识别优于统一滤波。问题四证明阶段归一化基线是实验集预测的核心，扰动残差只应作为局部修正。问题五给出了可执行的三级预警机制，并通过全组合变量筛选和逆速度法增强了工程可信度。

总体来看，本文模型具有三个特征：一是结果完整，覆盖题目所需全部表格、节点、预测值和预警阈值；二是证据充分，每个关键选择都有对照模型或敏感性分析支撑；三是工程可解释，所有数学模型均能对应边坡变形、入渗、孔压和速度预警机制。上述特征使本文方案不仅能完成竞赛计算任务，也具有推广到同类边坡监测场景的应用价值。
